% !PRINTED 339
\weberpage{339}
%<<CONT

1. \gesperrt{If $d\equiv 2, 3\pmod 4$, then all and only those
numbers of the field $\Omega$ are integral which are contained in the form}
\[
x+y\sqrt{d}
\]
\gesperrt{with rational integers $x, y$.}

But if $d\equiv 1\pmod 4$, then (3) is satisfied when $x$ and $y$
are both even or both odd, that is, when $x$ has the form
$2x_1-y$. If one then again replaces $x_1$ by $x$, it follows:

2. \gesperrt{If $d\equiv 1\pmod 4$, then all and only those numbers
in $\Omega$ are integral which are contained in the form}
\[
x+y\,\frac{-1+\sqrt{d}}{2}
\]
\gesperrt{with rational integers $x, y$.}

Thus in case 1. $(1,\sqrt{d})$, in case 2. $\left(1,\dfrac{-1+\sqrt{d}}{2}\right)$
is \gesperrt{a minimal basis of the field $\Omega$. The fundamental number or discriminant of the field $\Omega$} is therefore, in case 1.:
\weq{5}{\varDelta=\begin{vmatrix}1, & \sqrt{d}\\ 1, & -\sqrt{d}\end{vmatrix}^{2}=4d,}
in case 2.:
\weq{6}{\varDelta=\begin{vmatrix}1, & \dfrac{-1+\sqrt{d}}{2}\\[6pt] 1, & \dfrac{-1-\sqrt{d}}{2}\end{vmatrix}^{2}=d.}
In both cases one can therefore put the integer $\omega$ into the
form:
\weq{7}{\omega=\frac{x+y\sqrt{\varDelta}}{2},}
with the condition that
\weq{8}{4N(\omega)=x^2-\varDelta y^2\equiv 0\pmod 4}
shall hold.

The two forms of the minimal basis which we have obtained,
namely $(1,\sqrt{d})$ and $\left(1,\dfrac{-1+\sqrt{d}}{2}\right)$, are therefore:
\weq{9}{\begin{aligned}(1,0)&=\left(1,\tfrac{1}{2}\sqrt{\varDelta}\right), & \varDelta&\equiv 0,\\ &=\left(1,\frac{-1+\sqrt{\varDelta}}{2}\right), & \varDelta&\equiv 1.\end{aligned}\quad(\mathrm{mod}\ 4).}
Here the sign of $\sqrt{\varDelta}$ may be fixed in any manner whatever, but it must then be kept fixed throughout the same discussion. We will assume once and for all:
%CONT>>
